\paragraph{Observações.}
\begin{itemize}
  \item \textbf{FS inicial:} 2,093 nos três solos e nas três variantes (o estado inicial não depende do solo: sucção
    hidrostática com teto de 75~kPa e $\chi = 1$), contra 2,11 do artigo (Fig.~10) e 2,119 do relatório de outubro.
    O SRM de hoje é mais fino e dá FS um pouco menor.
  \item \textbf{Durante a chuva o FS cai sempre}, mais rápido quanto maior $I_r/k_s$. Na (iii), no f10 com 180~mm/h
    ($0{,}5k_s$) ele cruza 1 em 10,7~h e encosta num patamar de 0,84--0,85 a partir de 18~h: o FS do talude saturado
    até a superfície, o patamar inferior da Eq.~6 (0,81--0,82 no artigo). No f50 com 90~mm/h ($2{,}5k_s$) a queda
    acelera no fim da chuva (0,923 em 22~h, 0,832 em 24~h) e o FS fica abaixo de 1 de 19,3 a 32,3~h.
  \item \textbf{Depois da chuva o FS volta}, rápido no f10 (de 0,84 a 1,60 em 24~h), moderado no f50 (0,83 a 1,17) e
    devagar no f100 (1,55 a 1,62): a drenagem acompanha $k_s$. Com ar ativo a volta é mais rápida no f50 (1,39 em
    48~h com 36 e 90~mm/h, contra 1,27 e 1,17 na iii): o ar sai e $p_w$ cai.
  \item \textbf{Instante do mínimo:} 24~h (fim da chuva) em quase todos os casos. No f100 com 3,6 e 9~mm/h o FS fica
    num patamar depois da chuva; o mínimo (26--34~h) é 0,008 menor que o FS de 24~h, abaixo da resolução do SRM. Isto
    não contradiz a premissa ``o FS mínimo ocorre no fim da chuva'' (Hu 2024, Cho 2014) nestas chuvas.
\end{itemize}
